
\setcounter{footnote}{47}

\SGARefTarget{sga3:iv:section:5}\section*{5. Passage to the Quotient and Algebraic Structures}
\addcontentsline{toc}{section}{5. Passage to the Quotient and Algebraic Structures}

\SGARefTarget{sga3:iv:subsection:5:1}\subsection*{5.1. Principal Homogeneous Bundles}
\addcontentsline{toc}{subsection}{\SGARefLink{sga3:iv:subsection:5:1}{5.1}. Principal Homogeneous Bundles}

\noindent\SGARefTarget{sga3:iv:definition:5:1:0}\textbf{Definition 5.1.0.}\footnote{Editor's note: the numbering
\SGARefLink{sga3:iv:definition:5:1:0}{5.1.0} has been introduced for later reference.}
We recall (III \SGARefLink{sga3:iii:definition:0:1}{0.1}) that an object \(X\) equipped with a right action of
\(H\) is said to be \emph{formally principal homogeneous}\footnote{Editor's
note: one also says ``pseudo-torsor''; cf. EGA IV\textsubscript{4}, 16.5.15.
The more general notion of a formally homogeneous object (not necessarily
principal homogeneous) is defined in Addendum \SGARefLink{sga3:iv:definition:6:7:1}{6.7.1} at the end of this
expos\'e.} under \(H\) if the canonical morphism of functors
\[
  X\times H\longrightarrow X\times X,
  \qquad (x,h)\longmapsto(x,xh),
\]
is an isomorphism. Equivalently (cf. \emph{loc. cit.}), for every
\(S\in\operatorname{Ob}\mathscr C\), the set \(X(S)\) is formally principal
homogeneous under \(H(S)\): it is either empty or principal homogeneous
under \(H(S)\). In particular, when \(H\) acts on itself by right
translations, it is formally principal homogeneous under itself.

\medskip
\noindent\SGARefTarget{sga3:iv:definition:5:1:1}\textbf{Definition 5.1.1.}
An object \(X\) with an action of \(H\) is said to be \emph{trivial} if it
is isomorphic, as an \(H\)-object, to \(H\) with its translation action.

\medskip
\noindent\SGARefTarget{sga3:iv:proposition:5:1:2}\textbf{Proposition 5.1.2.}
\textit{If \(X\) is formally principal homogeneous under \(H\), there is an
isomorphism}
\[
  \Gamma(X)\xrightarrow{\sim}
  \operatorname{Isom}_{H\text{-obj.}}(H,X)
\]
\textit{of principal homogeneous sets under \(\Gamma(H)\).}

To a section \(x\) of \(X\), associate the morphism \(H\to X\) defined on
sets by \(h\mapsto xh\). \SGARefLink{sga3:iv:proposition:5:1:2}{The assertion} follows immediately by reduction to
the set-theoretic case.

\medskip
\noindent\SGARefTarget{sga3:iv:corollary:5:1:3}\textbf{Corollary 5.1.3.}
\textit{There is an isomorphism of \(H\)-objects}
\[
  X\xrightarrow{\sim}\operatorname{Isom}_{H\text{-obj.}}(H,X).
\]

\medskip
\noindent\SGARefTarget{sga3:iv:corollary:5:1:4}\textbf{Corollary 5.1.4.}
\textit{An object with a group action is trivial if and only if it is
formally principal homogeneous and has a section.}

\medskip
\noindent\SGARefTarget{sga3:iv:definition:5:1:5}\textbf{Definition 5.1.5.}
Let \(\mathscr C\) be a category equipped with a topology. An \(S\)-object
\(X\) with an action of the \(S\)-group \(H\) is called a
\emph{principal homogeneous bundle under \(H\)} if it is locally trivial,
that is, if the following equivalent conditions hold:\footnote{Editor's
note: in this case one also says that \(X\) is an \(H\)-torsor.}
\begin{enumerate}
  \item[(i)]\SGARefTarget{sga3:iv:definition:5:1:5:item:i:01} the collection of morphisms \(T\to S\) for which the functor
  \(X\times_ST\) is trivial under \(H\times_ST\) is a refinement of \(S\);
  \item[(ii)]\SGARefTarget{sga3:iv:definition:5:1:5:item:ii:01} for a pretopology defining the given topology, there is a
  covering family \(\{S_i\to S\}\) such that, for every \(i\), the
  \(S_i\)-functor \(X\times_SS_i\) is trivial under the \(S_i\)-group
  functor \(H\times_SS_i\), equivalently, it has a section over \(S_i\).
\end{enumerate}


\medskip
\noindent\SGARefTarget{sga3:iv:proposition:5:1:6}\textbf{Proposition 5.1.6.}
\textit{Let \(\mathscr C\) be a category equipped with a topology
\(\mathcal T\), and let \((M)\) be a family of morphisms of \(\mathscr C\)
satisfying axioms \SGARefLink{sga3:iv:para:4:6:3:item:a:01}{(a)}--\SGARefLink{sga3:iv:para:4:6:3:item:e:mathcal:t:01}{\((e_{\mathcal T})\)} of \SGARefLink{sga3:iv:para:4:6:3}{4.6.3}. Let \(H\) be an
\(S\)-group whose structure morphism \(H\to S\) belongs to \((M)\), and let
\(P\) be an \(S\)-object with an action of \(H\). The following conditions
are equivalent:}
\begin{enumerate}
  \item[(i)]\SGARefTarget{sga3:iv:proposition:5:1:6:item:i:01} \textit{\(P\) is a principal homogeneous bundle under \(H\)
  (Definition \SGARefLink{sga3:iv:definition:5:1:5}{5.1.5});}
  \item[(ii)]\SGARefTarget{sga3:iv:proposition:5:1:6:item:ii:01} \textit{\(P\) is formally principal homogeneous under \(H\),
  and its structure morphism \(P\to S\) belongs to \((M)\);}
  \item[(iii)]\SGARefTarget{sga3:iv:proposition:5:1:6:item:iii:01} \textit{there exists a morphism \(S'\to S\) belonging to
  \((M)\) such that, after base change from \(S\) to \(S'\), \(P\) becomes
  trivial; that is, \(P\times_SS'\) is trivial under \(H\times_SS'\);}
  \item[(iv)]\SGARefTarget{sga3:iv:proposition:5:1:6:item:iv:01} \textit{\(H\) acts freely and \((M)\)-effectively on \(P\),
  and the quotient \(P/H\) is isomorphic to \(S\).}
\end{enumerate}

First, \SGARefLink{sga3:iv:proposition:5:1:6:item:ii:01}{(ii)} and \SGARefLink{sga3:iv:proposition:5:1:6:item:iv:01}{(iv)} are equivalent. In either case \(P\to S\) belongs to
\((M)\), hence is squarable, which ensures that the fiber products
\(H\times_SP\) and \(P\times_SP\) are representable. Condition \SGARefLink{sga3:iv:proposition:5:1:6:item:ii:01}{(ii)} plainly
implies \SGARefLink{sga3:iv:proposition:5:1:6:item:iii:01}{(iii)}: one may take \(S'=P\), since formal principal homogeneity
implies that \(P\times_SP\) is trivial under \(H\times_SP\) by \SGARefLink{sga3:iv:corollary:5:1:4}{5.1.4}, its
diagonal section providing a section. Condition \SGARefLink{sga3:iv:proposition:5:1:6:item:iii:01}{(iii)} plainly implies \SGARefLink{sga3:iv:proposition:5:1:6:item:i:01}{(i)},
because \(\{S'\to S\}\) is a covering family by axiom
\SGARefLink{sga3:iv:para:4:6:3:item:d:mathcal:t:01}{\((d_{\mathcal T})\)}.

It remains to prove \SGARefLink{sga3:iv:proposition:5:1:6:item:i:01}{(i)}\(\Rightarrow\)\SGARefLink{sga3:iv:proposition:5:1:6:item:ii:01}{(ii)}. The morphism of sheaves
\[
  P\times_SH\longrightarrow P\times_SP
\]
is locally an isomorphism and hence is an isomorphism (\SGARefLink{sga3:iv:corollary:4:5:8}{4.5.8}); thus \(P\) is
formally principal homogeneous. The structure morphism \(P\to S\) is
locally isomorphic to the structure morphism \(H\to S\), which belongs to
\((M)\). It therefore belongs to \((M)\) by axiom
\SGARefLink{sga3:iv:para:4:6:3:item:e:mathcal:t:01}{\((e_{\mathcal T})\)}.

The equivalence between \SGARefLink{sga3:iv:proposition:5:1:6:item:i:01}{(i)} and \SGARefLink{sga3:iv:proposition:5:1:6:item:iv:01}{(iv)} admits the following generalization.

\medskip
\noindent\SGARefTarget{sga3:iv:proposition:5:1:7}\textbf{Proposition 5.1.7.}
\textit{Under the same assumptions on \(\mathscr C\) and \((M)\), let
\(H\) be an \(S\)-group and let \(X\) be an \(S\)-object on which \(H\)
acts on the right. Suppose that \(H\to S\) belongs to \((M)\). The following
conditions are equivalent:}
\begin{enumerate}
  \item[(i)]\SGARefTarget{sga3:iv:proposition:5:1:7:item:i:01} \textit{\(H\) acts freely and \((M)\)-effectively on \(X\);}
  \item[(ii)]\SGARefTarget{sga3:iv:proposition:5:1:7:item:ii:01} \textit{there exists an \(S\)-morphism \(p:X\to Y\),
  compatible with the equivalence relation on \(X\) defined by the action
  of \(H\), such that the induced action of \(H\times_SY\) on \(X\) over
  \(Y\) makes \(X\) a principal homogeneous bundle under \(H_Y\) over
  \(Y\).}
\end{enumerate}
\textit{Under these conditions, \(p\) identifies \(Y\) with the quotient
\(X/H\).}

If \(p:X\to Y\) is compatible with the action of \(H\), then the induced
action of \(H\times_SY\) on \(X\) over \(Y\) defines the same equivalence
relation on \(X\) as the original action of \(H\), by the isomorphism
\[
  H_Y\times_YX\xrightarrow{\sim}H\times_SX.
\]
\SGARefLink{sga3:iv:proposition:5:1:7}{The proposition} follows from this observation and the preceding equivalence
\SGARefLink{sga3:iv:proposition:5:1:6:item:iv:01}{(iv)}\(\Leftrightarrow\)\SGARefLink{sga3:iv:proposition:5:1:6:item:i:01}{(i)}.

\medskip
\noindent\SGARefTarget{sga3:iv:corollary:5:1:7:1}\textbf{Corollary 5.1.7.1.}\footnote{Editor's note: this special
case has been added as a corollary; it will be used several times in the
following expos\'es.}
\textit{Let \(\mathscr C\) be a category having a final object and fiber
products, and let it be equipped with a topology \(\mathcal T\) coarser than
the canonical topology. Let \(f:G\to H\) be a morphism of
\(\mathscr C\)-groups, put \(K=\operatorname{Ker}(f)\), and suppose that
\(f\) is covering for \(\mathcal T\). Then \(H\) represents the sheaf
quotient \(G/K\), and \(f\) is a \(K_H\)-torsor. Equivalently, one may say
that \(G\) is a \(K\)-torsor over \(H\).}

Indeed, because \(f\) is covering, it is a universal effective epimorphism
(\SGARefLink{sga3:iv:proposition:4:4:3}{4.4.3}). Hence, by \SGARefLink{sga3:iv:scholium:3:3:3:1}{3.3.3.1}, \(H\) is the quotient of \(G\) by the
equivalence relation
\[
  \mathscr R(f)=G\times_HG.
\]
Moreover, the morphism
\[
  G\times K\longrightarrow G\times_HG,
  \qquad (g,k)\longmapsto(g,gk),
\]
is an isomorphism of \(K_G=G\times_HK_H\)-objects; its inverse is
\((g,g')\mapsto(g,g^{-1}g')\). Thus \(\mathscr R(f)\) is the equivalence
relation defined by \(K\). Since \(f:G\to H\) is covering, it is a
\(K_H\)-torsor by \SGARefLink{sga3:iv:proposition:5:1:6:item:ii:01}{5.1.6(ii)}, or directly by Definition \SGARefLink{sga3:iv:definition:5:1:5:item:ii:01}{5.1.5(ii)}.

We can now sharpen Theorem \SGARefLink{sga3:iv:theorem:4:6:6}{4.6.6} in the case of a quotient by a group action.

\medskip
\noindent\SGARefTarget{sga3:iv:proposition:5:1:8}\textbf{Proposition 5.1.8.}
\textit{Under the assumptions of \SGARefLink{sga3:iv:proposition:5:1:7}{5.1.7}, let \(F_0\) denote the functor over
\(S\) defined as follows: for each \(S'\to S\), \(F_0(S')\) is the set of
representable \(S'\)-subfunctors \(Z\) of \(X\times_SS'\) that are stable
under \(H\times_SS'\) and are principal homogeneous bundles under this
\(S'\)-group for the induced action (\SGARefLink{sga3:iv:lemma:3:2:2}{3.2.2}).}
\begin{enumerate}
  \item[(i)]\SGARefTarget{sga3:iv:proposition:5:1:8:item:i:01} \textit{The following conditions are equivalent:}
  \begin{enumerate}
    \item[(a)]\SGARefTarget{sga3:iv:proposition:5:1:8:item:a:01} \textit{the action of \(H\) on \(X\) is
    \((M)\)-effective and free;}\footnote{Editor's note: the words ``and
    free'' have been added.}
    \item[(b)]\SGARefTarget{sga3:iv:proposition:5:1:8:item:b:01} \textit{\(F_0\) is representable.}
  \end{enumerate}
  \textit{Under these conditions, \(F_0=X/H\).}
  \item[(ii)]\SGARefTarget{sga3:iv:proposition:5:1:8:item:ii:01} \textit{Let \((N)\) be a family of morphisms stable under
  base change, such that for every covering family
  \(\{S_i'\to S'\}\) and every family of morphisms
  \(\{T_i\to S_i'\}\) in \((N)\), every descent datum on the \(T_i\)
  relative to \(\{S_i'\to S'\}\) is effective. Suppose that
  \(X\times_SH\to X\times_SX\) belongs to \((N)\), and that \(X\) is
  squarable. Then the element \(p\in\operatorname{Hom}(X,F_0)\)
  corresponding to the subobject \(X\times_SH\) of \(X\times_SX\)
  identifies \(F_0\) with the sheaf quotient \(X/H\).}
\end{enumerate}

\SGARefTarget{sga3:iv:subsection:5:2}\subsection*{5.2. Group Structures and Passage to the Quotient}
\addcontentsline{toc}{subsection}{\SGARefLink{sga3:iv:subsection:5:2}{5.2}. Group Structures and Passage to the Quotient}

In this subsection we are concerned with the algebraic structures that can
be placed on the quotient (G/H) of a group by a subgroup. We first work in
the category of sheaves on \(\mathscr C\) for an arbitrary topology. Taking
the canonical topology and using \SGARefLink{sga3:iv:lemma:4:5:12}{4.5.12} will then give results concerning
universal effective passage to the quotient in \(\mathscr C\).

\medskip
\noindent\SGARefTarget{sga3:iv:proposition:5:2:1}\textbf{Proposition 5.2.1.}
\textit{Let \(u:H\to G\) be a monomorphism of sheaves of groups. There is a
unique \(G\)-object structure on the quotient sheaf \(G/H\) for which the
canonical morphism
\[
  p:G\longrightarrow G/H
\]
is a morphism of \(G\)-objects. This structure is functorial in the pair
\((G,H)\): given a commutative diagram
\[
  \begin{array}{ccc}
    H & \longrightarrow & G \\
    \downarrow && \downarrow\scriptstyle f \\
    H' & \longrightarrow & G',
  \end{array}
\]
the morphism \(f:G/H\to G'/H'\) of \SGARefLink{sga3:iv:lemma:3:2:3}{3.2.3} is compatible with the morphism
\(f\) between the acting groups.}

Indeed, \(G/H\) is the sheaf associated with the presheaf
\[
  i(G)/i(H):S\longmapsto G(S)/H(S).
\]
Since the functor \(a\) is left exact, it carries objects with a group
action to objects with a group action. The presheaf \(i(G)/i(H)\) is an
\(i(G)\)-object, and therefore
\[
  G/H=a\bigl(i(G)/i(H)\bigr)
\]
is an \(a(i(G))=G\)-object. This construction plainly has all the stated
properties.

\medskip
\noindent\SGARefTarget{sga3:iv:corollary:5:2:2}\textbf{Corollary 5.2.2.}
\textit{Let \(u:H\to G\) be a monomorphism of \(\mathscr C\)-groups, and
suppose that the action of \(H\) on \(G\) is universally effective. There
is a unique \(G\)-object structure on the quotient object
\(G/H\in\operatorname{Ob}\mathscr C\) for which \(p:G\to G/H\) is a
morphism of objects with an action. In the sense of \SGARefLink{sga3:iv:proposition:5:2:1}{the preceding
proposition}, this structure is functorial in the pair \((H,G)\), with
\(H\) acting universally effectively on \(G\).}

\medskip
\noindent\SGARefTarget{sga3:iv:proposition:5:2:3}\textbf{Proposition 5.2.3.}
\textit{Let \(u:H\to G\) be a monomorphism of sheaves of groups that
identifies \(H\) with an invariant subgroup sheaf of \(G\). There is a
unique sheaf-of-groups structure on the quotient sheaf \(G/H\) for which
the canonical morphism \(p:G\to G/H\) is a group morphism. This structure
is functorial in the pair \((H,G)\), with \(H\) invariant.}

The proof is similar to that of \SGARefLink{sga3:iv:proposition:5:2:1}{5.2.1}.

\medskip
\noindent\SGARefTarget{sga3:iv:corollary:5:2:4}\textbf{Corollary 5.2.4.}
\textit{Let \(u:H\to G\) be a monomorphism of \(\mathscr C\)-groups that
identifies \(H\) with an invariant subgroup of \(G\). Suppose that the
action of \(H\) on \(G\) is universally effective. There is a unique group
structure on the quotient object \(G/H\in\operatorname{Ob}\mathscr C\) for
which the canonical morphism \(G\to G/H\) is a group morphism. This
structure is functorial in the pair \((H,G)\), with \(H\) invariant and
acting universally effectively.}

The group structure of \(G/H\) can be characterized more explicitly.

\medskip
\noindent\SGARefTarget{sga3:iv:proposition:5:2:5}\textbf{Proposition 5.2.5.}
\textit{Under the assumptions of \SGARefLink{sga3:iv:corollary:5:2:4}{5.2.4}, let \(K\) be a
\(\mathscr C\)-group and let \(f:G\to K\) be a morphism. The following
conditions are equivalent:}
\begin{enumerate}
  \item[(i)]\SGARefTarget{sga3:iv:proposition:5:2:5:item:i:01} \textit{\(f\) is a group morphism compatible with the
  equivalence relation defined by \(H\);}
  \item[(ii)]\SGARefTarget{sga3:iv:proposition:5:2:5:item:ii:01} \textit{\(f\) is a group morphism whose restriction
  \(H\to K\) is the trivial morphism;}
  \item[(iii)]\SGARefTarget{sga3:iv:proposition:5:2:5:item:iii:01} \textit{\(f\) factors through a group morphism
  \(G/H\to K\).}
\end{enumerate}
\textit{In particular, there is an isomorphism, functorial in the group
\(K\),}
\[
  \operatorname{Hom}_{\mathscr C\text{-gr.}}(G/H,K)
  \xrightarrow{\sim}
  \bigl\{f\in\operatorname{Hom}_{\mathscr C\text{-gr.}}(G,K)
  \mid f\circ u=e\bigr\}.
\]

The equivalence of \SGARefLink{sga3:iv:proposition:5:2:5:item:i:01}{(i)} and \SGARefLink{sga3:iv:proposition:5:2:5:item:ii:01}{(ii)} is proved set-theoretically, and
\SGARefLink{sga3:iv:proposition:5:2:5:item:iii:01}{(iii)}\(\Rightarrow\)\SGARefLink{sga3:iv:proposition:5:2:5:item:ii:01}{(ii)} is immediate. The equivalence of \SGARefLink{sga3:iv:proposition:5:2:5:item:iii:01}{(iii)} and \SGARefLink{sga3:iv:proposition:5:2:5:item:ii:01}{(ii)}
follows from the formula
\[
  \operatorname{Hom}(G/H,K)
  \simeq \operatorname{Hom}\bigl(i(G)/i(H),K\bigr)
\]
and the definition of the group structure on \(G/H\).

\medskip
\noindent\SGARefTarget{sga3:iv:remark:5:2:6}\textbf{Remark 5.2.6.}
In the preceding situation, if the kernel of \(f\) is exactly \(H\), then
the morphism \(G/H\to K\) through which \(f\) factors is a monomorphism.
This follows at once from \SGARefLink{sga3:iv:proposition:3:3:4}{3.3.4}.

For sheaves of groups, \SGARefLink{sga3:iv:proposition:4:4:11}{4.4.11} can be sharpened as follows.

\medskip
\noindent\SGARefTarget{sga3:iv:proposition:5:2:7}\textbf{Proposition 5.2.7.}
\textit{Let \(G\) be a sheaf of groups and \(H\) an invariant subgroup
sheaf. For every subgroup sheaf \(K\) of \(G\) containing \(H\), let
\(K'=K/H\), regarded as a subgroup of \(G'=G/H\). Then}
\[
  K=K'\times_{G'}G,
\]
\textit{and the maps}
\[
  K\longmapsto K/H,
  \qquad
  K'\longmapsto K'\times_{G'}G
\]
\textit{give a bijection between the subgroup sheaves of \(G\) containing
\(H\) and the subgroup sheaves of \(G'\). Under this correspondence, the
invariant subgroup sheaves of \(G\) and of \(G'\) correspond.}

The first assertion follows readily from \SGARefLink{sga3:iv:proposition:4:4:11}{4.4.11} and \SGARefLink{sga3:iv:lemma:3:2:4}{3.2.4}. It remains to
show that \(K\) is invariant in \(G\) if and only if \(K'\) is invariant in
\(G'\). If \(K\) is invariant in \(G\), then the presheaf
\(i(K)/i(H)\) is invariant in \(i(G)/i(H)\); by the usual argument, the
same is true of the associated sheaves. Conversely, if \(K'\) is invariant
in \(G'\), then the fiber product \(K'\times_{G'}G\) is invariant in
\(G\), as is immediate.

Now let \(L\) be an arbitrary subgroup sheaf of \(G\). Denote by
\(\overline L\) the saturation of \(L\) for the equivalence relation
defined by \(H\); we also write \(\overline L=L\cdot H\).

\medskip
\noindent\SGARefTarget{sga3:iv:proposition:5:2:8}\textbf{Proposition 5.2.8.}
\textit{Under the preceding assumptions, \(L\cdot H\) is a subgroup sheaf
of \(G\) containing \(H\), and the image of \(L\) in \(G/H\) is
identified with}
\[
  (L\cdot H)/H\simeq L/(H\cap L).
\]

Indeed, let \(L'\) denote the image sheaf of \(L\) in \(G/H\). Under the
correspondence of \SGARefLink{sga3:iv:proposition:5:2:8}{the preceding proposition}, it is the subgroup sheaf of
\(G/H\) corresponding to \(L\cdot H\). Since \(L\to L'\) is covering,
and hence a universal effective epimorphism of sheaves, it follows from
\SGARefLink{sga3:iv:proposition:4:4:9}{4.4.9} that \(L'\) is identified with the quotient of \(L\) by the kernel
of \(L\to L'\), which is plainly \(H\cap L\).

Finally, consider a sheaf of groups \(G\), a subgroup sheaf \(K\), and an
invariant subgroup sheaf \(H\) of \(K\). We first define a right action of
the sheaf of groups \(H\backslash K\;(=K/H)\) on \(G/H\). The group
\(K\) acts on \(G\) by right translations. Since \(H\) is invariant in
\(K\), this action is compatible with the equivalence relation defined by
the action of \(H\), and therefore induces an action of \(K\) on \(G/H\),
that is, a morphism
\[
  K^{\circ}\longrightarrow\operatorname{Aut}(G/H)
\]
from the group opposite to \(K\). Since the target is a sheaf (\SGARefLink{sga3:iv:corollary:4:5:13}{4.5.13}) and
this morphism is trivial on \(H\), it factors through \(K/H\) and defines
the desired action. The right and left actions of \(G\) on itself commute;
hence the actions of \(G\) and \(K/H\) on \(G/H\) commute.

\medskip
\noindent\SGARefTarget{sga3:iv:proposition:5:2:9}\textbf{Proposition 5.2.9.}
\textit{Under the preceding assumptions, \(K/H\) acts freely on the right
on \(G/H\), and there is a canonical isomorphism of \(G\)-objects}
\[
  (G/H)/(K/H)\simeq G/K.
\]
\textit{If \(K\) is invariant in \(G\), in which case \(K/H\) is invariant
in \(G/H\) by \SGARefLink{sga3:iv:proposition:5:2:7}{5.2.7}, this isomorphism respects the group structures on
both sides.}

There is an isomorphism of presheaves
\[
  \bigl(i(G)/i(H)\bigr)\big/\bigl(i(K)/i(H)\bigr)
  \xleftarrow{\sim} i(G)/i(K),
\]
which respects the \(i(G)\)-object structures. Applying the functor \(a\)
to this relation gives the asserted result.

\medskip
\noindent\SGARefTarget{sga3:iv:corollary:5:2:10}\textbf{Corollary 5.2.10.}
\textit{Let \(G\) be a \(\mathscr C\)-group, \(K\) a
\(\mathscr C\)-subgroup of \(G\), and \(H\) an invariant
\(\mathscr C\)-subgroup of \(K\). Let \((M)\) be a family of morphisms of
\(\mathscr C\) satisfying axioms \SGARefLink{sga3:iv:para:4:6:3:item:a:01}{(a)}--\SGARefLink{sga3:iv:para:4:6:3:item:e:mathcal:t:01}{\((e_{\mathcal T})\)}. Suppose that
the right action of \(H\) on \(G\) (respectively, on \(K\)) is
\((M)\)-effective. Then \(K/H\) acts naturally and freely on the right on
\(G/H\), and this action commutes with that of \(G\). The following
conditions are equivalent:}
\begin{enumerate}
  \item[(i)]\SGARefTarget{sga3:iv:corollary:5:2:10:item:i:01} \textit{the action of \(K\) on \(G\) is
  \((M)\)-effective;}
  \item[(ii)]\SGARefTarget{sga3:iv:corollary:5:2:10:item:ii:01} \textit{the action of \(K/H\) on \(G/H\) is
  \((M)\)-effective.}
\end{enumerate}
\textit{Under these conditions, there is an isomorphism of
objects\footnote{Editor's note: of \(\mathscr C\).} with a \(G\)-action}
\[
  (G/H)/(K/H)\simeq G/K.
\]

\SGARefTarget{sga3:iv:subsection:5:3}\subsection*{5.3. Use of Effectivity Criteria: Noether's Theorem}
\addcontentsline{toc}{subsection}{\SGARefLink{sga3:iv:subsection:5:3}{5.3}. Use of Effectivity Criteria: Noether's Theorem}
\footnote{Editor's note: in fact, the isomorphisms established in
\SGARefLink{sga3:iv:proposition:5:2:8}{5.2.8}--\SGARefLink{sga3:iv:corollary:5:2:10}{5.2.10} deserve to be called ``Noether isomorphism theorems.''}
Let \(\mathscr C\), \(\mathcal T\), and \((M)\) be as usual. Let \((N)\)
be a family of morphisms satisfying axioms \SGARefLink{sga3:iv:subsection:4:7:item:a:01}{(a)} and \SGARefLink{sga3:iv:subsection:4:7:item:f:m:01}{\((f_M)\)} of \SGARefLink{sga3:iv:subsection:4:7}{4.7}.
Combining \SGARefLink{sga3:iv:proposition:5:2:7}{5.2.7} and \SGARefLink{sga3:iv:proposition:4:7:2}{4.7.2} gives \SGARefLink{sga3:iv:proposition:5:3:1}{the following result}.

\medskip
\noindent\SGARefTarget{sga3:iv:proposition:5:3:1}\textbf{Proposition 5.3.1.}
\textit{Let \(G\) be a \(\mathscr C\)-group, and let \(H\) be an invariant
\(\mathscr C\)-subgroup of \(G\) that acts \((M)\)-effectively on \(G\).
For every \(\mathscr C\)-subgroup \(K\) of \(G\) containing \(H\) such
that \(K\to G\) belongs to \((N)\), the action of \(H\) on \(K\) is
\((M)\)-effective, and the quotient \(K/H=K'\) is a
\(\mathscr C\)-subgroup of \(G/H=G'\) such that \(K'\to G'\) belongs to
\((N)\). The map \(K\mapsto K'\) is a bijection between the set of
\(\mathscr C\)-subgroups \(K\) of \(G\) that contain \(H\) and for which
\(K\to G\) belongs to \((N)\), and the set of \(\mathscr C\)-subgroups
\(K'\) of \(G'\) for which \(K'\to G'\) belongs to \((N)\). Its inverse
is}
\[
  K'\longmapsto K\times_{G'}G.
\]
\textit{Under this correspondence, the invariant subgroups of \(G\) and
of \(G'\) correspond.}

\medskip
\noindent\SGARefTarget{sga3:iv:corollary:5:3:2}\textbf{Corollary 5.3.2.}
\textit{If \(H\to G\) belongs to \((N)\), then \(\mathscr C\) has a final
object \(e\), and the unit section \(e\to G/H\) belongs to
\((N)\).}\footnote{Editor's note: this follows by applying \SGARefLink{sga3:iv:proposition:5:3:1}{the proposition}
with \(K=H\).}
